\documentclass[11pt,a4paper]{article}

\usepackage[margin=2.5cm]{geometry}
\usepackage{amsmath,amssymb,bm}
\usepackage{array}
\usepackage{booktabs}
\usepackage{graphicx}
\usepackage{hyperref}
\usepackage{microtype}
\usepackage{siunitx}

\hypersetup{colorlinks=true, linkcolor=blue, citecolor=blue, urlcolor=blue}
\sisetup{per-mode=symbol}
\newcommand{\chfour}{\ensuremath{\mathrm{CH_4}}}
\newcommand{\cv}{\ensuremath{C_V}}
\newcommand{\cpheat}{\ensuremath{C_P}}

\title{Efficient Heat Capacity of Methane-Loaded MOF-5 with Classical Molecular Dynamics and Machine-Learned Hessians}
\author{\textit{Author name}}
\date{\today}

\begin{document}

\maketitle

\begin{abstract}
This report develops an efficient heat-capacity workflow for methane-loaded
MOF-5 using a machine-learned interatomic potential (MLIP). Classical
flexible-cell NPT molecular dynamics (MD) samples the loaded system and retains
anharmonic methane and host--guest motion. Automatic differentiation (AD)
through the matching PET-JAX model provides sparse Cartesian Hessians and
normal-mode frequencies at optimized local minima. Following the approximation
tested by Kapil \emph{et al.}, the classical anharmonic heat capacity is
corrected by the difference between quantum and classical harmonic heat
capacities of the same loaded composition. The already-equilibrated empty
MOF-5 structure is relaxed directly and evaluated with one reference Hessian;
no empty-framework MD campaign is required. The workflow is designed to retain
the main quantum correction without the cost of full-system path-integral MD,
while recognizing that it is an approximation whose low-temperature and
model-dependent accuracy must be validated.
\end{abstract}

\section{Scientific motivation and scope}

Metal--organic frameworks (MOFs) combine inorganic nodes with organic linkers
to form porous crystalline networks. Their vibrational spectrum contains
high-frequency local bond vibrations as well as soft linker, pore, and
collective framework modes. Because the contribution of a quantum vibrational
mode depends strongly on the ratio of its energy to $k_{\mathrm B}T$, a broad
frequency distribution produces a strongly temperature-dependent heat
capacity. Guest adsorption adds translational, rotational, and host--guest
degrees of freedom that may be weakly confined and strongly anharmonic.

The material considered here is MOF-5. Its conventional cubic cell contains
424 framework atoms with composition
$\mathrm{C}_{192}\mathrm{H}_{96}\mathrm{O}_{104}\mathrm{Zn}_{32}$. In the
loaded calculations, 100 methane molecules add 500 atoms, giving a 924-atom
system with composition
$\mathrm{C}_{292}\mathrm{H}_{496}\mathrm{O}_{104}\mathrm{Zn}_{32}$.
The principal objectives are:

\begin{enumerate}
  \item to sample the anharmonic thermodynamics of methane-loaded MOF-5 with
        classical NPT MD and independent replicas;
  \item to compute force constants efficiently by sparse AD at optimized empty
        and loaded structures;
  \item to obtain and converge the loaded-system harmonic quantum correction;
  \item to combine that correction with the classical enthalpy derivative to
        estimate the loaded-system \cpheat{} efficiently; and
  \item to establish a reproducible route from short cluster validation jobs
        to statistically meaningful production calculations.
\end{enumerate}

The workflow has three computational stages. Loaded-system MD supplies the
classical anharmonic enthalpy curve. Fixed-cell optimizations followed by AD
Hessians supply quantum-minus-classical harmonic corrections. A final analysis
combines the two terms with consistent composition and mass normalization. The
empty-framework Hessian is retained as a reference and validation calculation,
not as a replacement for the loaded-system correction.

\section{Theory}
\label{sec:theory}

\subsection{Machine-learned interatomic potentials}

\subsubsection{Learning the potential-energy surface}

Within the Born--Oppenheimer approximation, the electronic ground-state
energy defines a potential-energy surface (PES) for the nuclei. An
interatomic potential replaces an electronic-structure calculation by an
approximation to this map,
\begin{equation}
  (\bm R,\bm h,\bm Z)\longmapsto
  U(\bm R,\bm h,\bm Z),
\end{equation}
where $\bm R=(\bm R_1,\ldots,\bm R_N)$ are the nuclear coordinates,
$\bm h$ is the periodic cell, and $\bm Z$ lists the chemical species. A
machine-learned interatomic potential (MLIP) represents this map by a
differentiable function with parameters $\theta$, fitted to energies and
forces from electronic-structure calculations. Modern local MLIPs usually
write the extensive energy as
\begin{equation}
  U_\theta(\bm R,\bm h,\bm Z)
  = U_{\mathrm{base}}(\bm Z)
  + \sum_{i=1}^{N}\varepsilon_i(\mathcal N_i;\theta),
  \label{eq:mlip-energy}
\end{equation}
where $U_{\mathrm{base}}$ is an optional composition-dependent baseline and
$\mathcal N_i$ is the learned environment of atom $i$. The decomposition
makes the cost approximately linear in $N$ at fixed density and makes the
energy size extensive. The individual atomic energies are a model
construction rather than uniquely defined physical observables.

The force field is obtained by differentiating the scalar energy,
\begin{equation}
  \bm F_i=-\frac{\partial U_\theta}{\partial \bm R_i},
  \label{eq:mlip-force}
\end{equation}
so the predicted forces are conservative, provided the same differentiable
energy model is used throughout. Cell derivatives similarly give the virial
or stress when the model and its interface support them. Automatic
differentiation can be applied a second time to obtain the Cartesian Hessian
used later in this project.

\subsubsection{From an atomic structure to a graph}

A graph neural network (GNN) represents the atoms by vertices and nearby
interactions by directed edges. For a periodic system, an edge
$(j,\bm n)\to i$ is included when
\begin{align}
  \bm r_{ij\bm n}
    &= \bm R_j+\bm h\bm n-\bm R_i,
  & r_{ij\bm n}&=\lVert\bm r_{ij\bm n}\rVert < r_{\mathrm c},
  \label{eq:periodic-edge}
\end{align}
where $\bm n\in\mathbb Z^3$ identifies a periodic image and
$r_{\mathrm c}$ is a cutoff. This construction removes dependence on the
absolute origin. A sufficiently smooth cutoff function $f_{\mathrm c}(r)$,
whose value and relevant derivatives vanish at $r_{\mathrm c}$, suppresses a
neighbor's contribution as it leaves the environment. Smooth forces and
Hessians require consistent cutoff treatment throughout the network,
including attention normalization and energy readout, rather than continuity
of the cutoff value alone.

Neural networks operate on fixed-width numerical representations. Chemical
identity is therefore mapped from a discrete species label to a trainable
embedding vector,
\begin{equation}
  \bm h_i^{(0)}=\bm E_Z[Z_i]\in\mathbb R^d,
  \label{eq:species-embedding}
\end{equation}
where each row of the embedding matrix $\bm E_Z$ is learned during training.
Geometry may be embedded through radial basis functions,
\begin{equation}
  \bm e_{ij}^{\mathrm{rad}}
  =f_{\mathrm c}(r_{ij})
   \bigl[\phi_1(r_{ij}),\ldots,\phi_K(r_{ij})\bigr],
  \label{eq:radial-embedding}
\end{equation}
and, in models that retain directional information, by the displacement
$\bm r_{ij}$, its unit vector, spherical harmonics, or learned functions of
these quantities. Here ``embedding vector'' means a vector in latent feature
space; it need not transform as a physical three-dimensional vector.

\subsubsection{Message passing and many-body representations}

At message-passing layer $\ell$, a generic GNN constructs a message on each
edge, aggregates the incoming messages in a permutation-independent way, and
updates the central representation:
\begin{align}
  \bm m_{ij}^{(\ell)}
    &=M_\ell\bigl(\bm h_i^{(\ell)},\bm h_j^{(\ell)},
                   \bm e_{ij}\bigr),
  \\
  \bm m_i^{(\ell)}
    &=\sum_{j\in\mathcal N_i}\bm m_{ij}^{(\ell)},
  \\
  \bm h_i^{(\ell+1)}
    &=U_\ell\bigl(\bm h_i^{(\ell)},\bm m_i^{(\ell)}\bigr).
  \label{eq:message-passing}
\end{align}
The learned maps $M_\ell$ and $U_\ell$ are commonly multilayer
perceptrons or attention blocks. Although the graph initially contains only
pairwise edges, repeated updates create many-body features: after $L$ layers,
an atom can depend on paths extending up to $L$ graph hops. A readout map then
converts the final features into the local energies in
Eq.~\eqref{eq:mlip-energy}.

\subsubsection{Point Edge Transformer architecture}
\label{sec:pet-architecture}

The Point Edge Transformer (PET) introduced by Pozdnyakov and Ceriotti
\cite{pet} is the architecture underlying the later PET-MAD potential
\cite{petmad}. PET preserves a distinct latent message on each directed edge.
For a central atom $i$, its transformer processes the incoming edge messages
together and produces a separate outgoing message for each neighbor. This
avoids compressing the entire neighborhood into a single atom representation
before forming outgoing messages. Attention still performs weighted sums
internally; the distinction is that the edge identities survive the update.

In the original formulation (Ref.~\cite{pet}, Appendix A), neighbor species,
relative Cartesian position, and an incoming message
$\bm m_{j\to i}^{(\ell-1)}\in\mathbb R^{d_{\mathrm{PET}}}$ are encoded as
\begin{align}
  \bm p_{ij}^{(\ell)}
    &=\operatorname{SiLU}\bigl(\bm W_r^{(\ell)}\bm r_{ij}
                                      +\bm b_r^{(\ell)}\bigr),
  \\
  \bm t_{ij}^{(\ell)}
    &=\mathrm{MLP}_\ell\!\left(
      \bm m_{j\to i}^{(\ell-1)}\Vert
      \bm E_n^{(\ell)}[Z_j]\Vert
      \bm p_{ij}^{(\ell)}\right),
  \label{eq:pet-token}
\end{align}
where $\Vert$ denotes concatenation. Each ingredient has width
$d_{\mathrm{PET}}$, and the MLP compresses their concatenation from
$3d_{\mathrm{PET}}$ to $d_{\mathrm{PET}}$. In the first block the incoming
message is absent, giving $2d_{\mathrm{PET}}$ input features. A separate
central token, $\bm t_{i0}^{(\ell)}=\bm E_c^{(\ell)}[Z_i]$, uses a distinct
species embedding. These are latent feature vectors with no prescribed
rotation law. Later PET implementations can modify the encoding and feature
dimensions; the project-specific settings are given in
Section~\ref{sec:pet-implementation}.

A transformer mixes the central and neighbor tokens within one environment.
For one attention head, with $a,b$ indexing these tokens,
\begin{align}
  \bm q_a&=\bm W_Q\bm t_a,
  &\bm k_b&=\bm W_K\bm t_b,
  &\bm v_b&=\bm W_V\bm t_b,
  \\
  \alpha_{ab}
    &=\frac{\exp(\bm q_a^{\mathrm T}\bm k_b/\sqrt{d_k})
                    c_b}
            {\sum_c\exp(\bm q_a^{\mathrm T}\bm k_c/\sqrt{d_k})
                    c_c},
  &\bm y_a&=\sum_b\alpha_{ab}\bm v_b.
  \label{eq:self-attention}
\end{align}
where $c_0=1$ for the central token and $c_j=f_{\mathrm c}(r_{ij})$ for
neighbor $j$. The cutoff enters before normalization, so a departing neighbor
vanishes smoothly from both the weighted sum and its denominator.
Multi-head attention repeats this operation in learned subspaces, followed
by feed-forward layers and residual connections. No arbitrary neighbor order
enters the calculation, so permuting neighbors permutes their output tokens
consistently.

Attention makes each neighbor token depend on the other neighbors, allowing
many-body interactions within a single message-passing block. Adding
attention layers increases local expressive power without extending that
block's neighborhood. Adding message-passing blocks communicates information
between environments and enlarges the receptive field. These are separate
architecture choices: the water benchmark in Ref.~\cite{pet} shows that
deeper local transformers can outperform a larger number of shallow
message-passing blocks at fixed total attention-layer count. Local attention
is quadratic in the number of tokens per environment, while the total cost
remains approximately linear in atom count for a fixed architecture and
bounded neighborhood size.

The output neighbor tokens supply outgoing messages, updated with residual
connections to preceding messages. Each block also contributes to the energy.
For the summed-edge variant in the original paper,
\begin{equation}
  U_\theta=\sum_i b_{Z_i}
  +\sum_{\ell=1}^{L}\sum_i\left[
      H_c^{(\ell)}(\bm x_{i0}^{(\ell)})
      +\sum_{j\in\mathcal N_i}f_{\mathrm c}(r_{ij})
                         H_n^{(\ell)}(\bm x_{ij}^{(\ell)})\right],
  \label{eq:pet-readout}
\end{equation}
where $\bm x_{i0}$ and $\bm x_{ij}$ are central and neighbor output tokens,
$H_c$ and $H_n$ are learned readout heads, and $b_Z$ are fitted species
baselines. Appendix A of Ref.~\cite{pet} also describes a cutoff-weighted
average of edge predictions. The edge terms are many-body learned
contributions, not physical pair energies. Contributions from all blocks
enter the total, and cutoff weighting is required in the readout as well as
in attention.

\subsubsection{Physical symmetries}

For an isolated non-magnetic system without external fields, a valid scalar
energy should be unchanged by translation, rotation or reflection of the
whole structure and by permutation of identical atoms:
\begin{equation}
  U(\bm R)=U(\bm R+\bm a)
  =U(\bm Q\bm R)
  =U(\bm P\bm R),
  \qquad \bm Q\in O(3).
  \label{eq:energy-symmetry}
\end{equation}
An invariant scalar is unchanged by a symmetry operation, whereas an
equivariant vector changes in the corresponding way; for example,
$\bm F_i(\bm Q\bm R)=\bm Q\bm F_i(\bm R)$. Translation invariance follows
from using relative positions, and the shared local functions and symmetric
aggregation enforce permutation invariance. Many MLIPs impose rotational
symmetry exactly through invariant geometric descriptors or
$E(3)$-equivariant features. PET instead processes Cartesian displacement
vectors with an unconstrained transformer and learns rotational behavior from
random rotational augmentation \cite{pet,petmad}. Thus the PET-MAD backbone
is not intrinsically exactly rotation invariant. The PET-MAD publication
reports rotational discrepancies usually smaller than its prediction errors
\cite{petmad}; this is not a validation of the particular MOF-5 structures or
checkpoint used here. For periodic structures, rotation tests must rotate
both positions and cell vectors.

\subsubsection{Exact rotational symmetrization with ECSE}

The original PET paper also introduces the Equivariant Coordinate System
Ensemble (ECSE), a separate construction that imposes exact rotational
symmetry on a smooth, translation- and permutation-invariant backbone
\cite{pet}. Each ordered noncollinear neighbor pair $(j,k)$ defines a local
frame $\bm Q_{ijk}$ that rotates with the atomic environment. For a scalar
local prediction, the symmetrized model is
\begin{equation}
  \varepsilon_i^{\mathrm{ECSE}}=
    \frac{\sum_{j\ne k}w_{ijk}\,
          \varepsilon_i^0(\bm Q_{ijk}^{\mathrm T}\mathcal N_i)}
         {\sum_{j\ne k}w_{ijk}},
  \label{eq:ecse}
\end{equation}
where the notation indicates expressing all relevant displacement vectors in
that frame. Vector or tensor predictions are transformed back to the
original frame before averaging. The frames in the paper use proper
rotations, $SO(3)$; reflection invariance in Eq.~\eqref{eq:energy-symmetry}
is an additional requirement, not a consequence of this construction alone.

Radial weights suppress frames whose defining neighbors approach the cutoff,
and angular weights suppress nearly collinear pairs. An adaptive inner radius
and smooth pruning reduce the number of frames. Fully collinear environments
require a fallback; the paper demonstrates an auxiliary invariant model.
Using a smooth ensemble avoids the discontinuities caused by selecting only
the two nearest neighbors. This frame-selection radius is distinct from the
adaptive neighbor cutoff in the later PET implementation used here.

Forces and Hessians must differentiate the full expression in
Eq.~\eqref{eq:ecse}, including the position dependence of frames and weights.
They cannot be obtained simply by averaging derivatives of the backbone.
The paper shows that tight frame-selection parameters can amplify derivative
errors even when the energy remains mathematically smooth. Its
proof-of-principle ECSE implementation incurs roughly three orders of
magnitude of inference overhead (Ref.~\cite{pet}, Section 6 and Appendix F).
ECSE is not enabled in this project's workflow; its exact symmetry properties
must therefore be distinguished from those of the deployed PET backbone.

\subsubsection{Training and domain of validity}

Given reference configurations $s$, a typical energy--force training
objective is
\begin{equation}
  \mathcal L(\theta)=\sum_s\left[
    w_E\left(\frac{U_\theta^{(s)}-U_{\mathrm{ref}}^{(s)}}{N_s}\right)^2
    +\frac{w_F}{3N_s}\sum_{i\alpha}
      \left(F_{i\alpha,\theta}^{(s)}
                 -F_{i\alpha,\mathrm{ref}}^{(s)}\right)^2
  \right],
  \label{eq:mlip-loss}
\end{equation}
possibly augmented by stress terms and regularization. Because the force
labels are derivatives of the PES, each configuration supplies $3N_s$ local
slope constraints in addition to its energy. Optimization adjusts the
embeddings, message-passing transformations and readout weights to minimize
this loss. The original PET experiments first subtract fitted species
baselines and use a loss that rescales energy and force errors by moving
validation-set mean squared errors, together with Adam and random rotational
augmentation (Ref.~\cite{pet}, Appendix B). Equation~\eqref{eq:mlip-loss}
is a generic objective, not that paper's exact normalization. PET-MAD was
trained on a chemically diverse DFT dataset rather than
on MOF-5 alone, making it a universal pretrained potential \cite{petmad}.

An MLIP interpolates the reference PES represented by its training data; it
does not replace the approximations of the underlying electronic-structure
method and need not extrapolate reliably to unfamiliar bonding, charge states,
very short distances or thermodynamic phases. Energy and force test errors,
trajectory stability and uncertainty indicators are therefore necessary but
distinct checks. This project also uses second derivatives, which can be more
sensitive than energies or forces and were not directly constrained by the
usual loss in Eq.~\eqref{eq:mlip-loss}. Curvatures and vibrational spectra must
therefore be validated explicitly. The original PET benchmarks include
isolated distorted methane, water, molecular collisions, crystals, and
high-entropy alloys, but do not validate methane adsorption in MOF-5. The
paper also reports weaker high-temperature alloy extrapolation despite strong
hold-out accuracy, illustrating the distinction between expressiveness and
transferability \cite{pet}.

\subsection{Classical molecular dynamics}
\label{sec:classical-md-theory}

\subsubsection{Hamiltonian dynamics and statistical averages}

Molecular dynamics (MD) generates a trajectory through phase space by
propagating the nuclear positions and momenta on a chosen PES. For classical
point nuclei with masses $m_i$, the Hamiltonian is
\begin{equation}
  \mathcal H(\bm R,\bm P)
  =\sum_{i=1}^{N}\frac{\bm P_i^2}{2m_i}+U(\bm R),
  \label{eq:classical-hamiltonian}
\end{equation}
and Hamilton's equations give
\begin{equation}
  \dot{\bm R}_i=\frac{\bm P_i}{m_i},
  \qquad
  \dot{\bm P}_i=-\frac{\partial U}{\partial\bm R_i}
  =\bm F_i.
  \label{eq:hamilton-equations}
\end{equation}
In exact microcanonical dynamics, $\mathcal H$ is conserved. If the dynamics
is ergodic and the trajectory is sufficiently long, an equilibrium average
can be estimated by a time average,
\begin{equation}
  \langle A\rangle
  =\lim_{\tau\to\infty}\frac{1}{\tau}
   \int_0^\tau A\bigl(\bm R(t),\bm P(t)\bigr)\,\mathrm dt.
  \label{eq:ergodic-average}
\end{equation}
Finite trajectories contain correlated samples, so their statistical
precision is controlled by the simulation length relative to the integrated
autocorrelation time, not simply by the number of saved frames.

The continuous equations are replaced by a finite-timestep integrator. A
representative symplectic scheme is velocity Verlet,
\begin{align}
  \bm P_i^{n+1/2}
    &=\bm P_i^n+\frac{\Delta t}{2}\bm F_i(\bm R^n),
  \\
  \bm R_i^{n+1}
    &=\bm R_i^n+\Delta t\,\frac{\bm P_i^{n+1/2}}{m_i},
  \\
  \bm P_i^{n+1}
    &=\bm P_i^{n+1/2}
      +\frac{\Delta t}{2}\bm F_i(\bm R^{n+1}).
  \label{eq:velocity-verlet}
\end{align}
It is time reversible and symplectic, and has global error of order
$\mathcal O(\Delta t^2)$. It approximately conserves a nearby ``shadow''
Hamiltonian rather than the exact one. The timestep must resolve the fastest
vibration; an apparently controlled temperature does not by itself rule out
integration error.

\subsubsection{Temperature, ensembles and thermostats}

For $g$ unconstrained kinetic degrees of freedom, the instantaneous kinetic
temperature is conventionally defined through equipartition,
\begin{equation}
  T_{\mathrm{kin}}=\frac{2K}{gk_{\mathrm B}},
  \qquad
  K=\sum_i\frac{\bm P_i^2}{2m_i}.
  \label{eq:kinetic-temperature}
\end{equation}
Temperature is an ensemble property; $T_{\mathrm{kin}}$ fluctuates along a
finite canonical trajectory. Constant-energy (NVE) dynamics samples the
microcanonical ensemble. To sample the canonical ensemble, whose stationary
phase-space density is
\begin{equation}
  \rho_{NVT}(\bm R,\bm P)\propto
  \exp\bigl[-\beta\mathcal H(\bm R,\bm P)\bigr],
  \qquad \beta=(k_{\mathrm B}T)^{-1},
  \label{eq:canonical-density}
\end{equation}
the physical variables may be coupled to a thermostat.

For example, Langevin dynamics adds friction and random forces,
\begin{align}
  m_i\ddot{\bm R}_i
    &=\bm F_i-m_i\gamma\dot{\bm R}_i+\bm\eta_i(t),
  \\
  \left\langle\eta_{i\alpha}(t)\eta_{j\beta}(t')\right\rangle
    &=2m_i\gamma k_{\mathrm B}T\,
      \delta_{ij}\delta_{\alpha\beta}\delta(t-t').
  \label{eq:langevin}
\end{align}
The second relation is the fluctuation--dissipation condition that makes the
canonical distribution stationary. Weak friction perturbs the natural
dynamics less but exchanges heat slowly; strong friction controls temperature
rapidly but can overdamp motion. Nos\'e--Hoover-chain thermostats instead add
a sequence of deterministic extended variables and are designed to generate
the same canonical distribution while improving ergodicity over a single
Nos\'e--Hoover variable \cite{nhc}. Thermostat choice affects sampling
efficiency and dynamical correlation functions even when equilibrium averages
are unchanged.

An isothermal--isobaric simulation additionally evolves the volume or cell so
that the target density is proportional to
$\exp[-\beta(\mathcal H+PV)]$. Flexible-cell algorithms must include the
appropriate cell kinetic variables, Jacobian and virial terms; the
Martyna--Tobias--Klein scheme is one such extended-system formulation
\cite{mtk}. Reliable NPT dynamics therefore requires a potential with
validated cell derivatives or stresses. NVE, NVT and NPT trajectories answer
different thermodynamic questions and should not be interchanged solely
because their mean temperatures are similar.

\subsection{Path-integral molecular dynamics}
\label{sec:pimd-theory}

\subsubsection{Quantum statistical mechanics as a ring polymer}

Classical MD neglects the wave-like character of the nuclei. For
distinguishable quantum nuclei on a single Born--Oppenheimer surface, the
Hamiltonian operator and canonical partition function are
\begin{equation}
  \hat{\mathcal H}=\sum_i\frac{\hat{\bm P}_i^2}{2m_i}
  +U(\hat{\bm R}),
  \qquad
  Z=\operatorname{Tr}\left[e^{-\beta\hat{\mathcal H}}\right].
  \label{eq:quantum-partition}
\end{equation}
The kinetic and potential operators do not commute. Splitting the imaginary
time $\beta\hbar$ into $P$ slices with a Trotter factorization maps the
quantum partition function, in the limit $P\to\infty$, onto a classical
extended system. Its primitive ring-polymer Hamiltonian is
\begin{align}
  Z_P&\propto\int\mathrm d\bm R\,\mathrm d\bm P\,
     e^{-\beta_P H_P},
  &\beta_P&=\frac{\beta}{P},
  \\
  H_P&=\sum_{s=1}^{P}\left[
    \sum_i\frac{|\bm P_i^{(s)}|^2}{2m_i}
    +\sum_i\frac{m_i\omega_P^2}{2}
      |\bm R_i^{(s)}-\bm R_i^{(s+1)}|^2
    +U(\bm R^{(s)})\right],
  &\omega_P&=\frac{P}{\beta\hbar},
  \label{eq:ring-polymer-hamiltonian}
\end{align}
with cyclic boundary condition
$\bm R_i^{(P+1)}=\bm R_i^{(1)}$ \cite{ipi}. Each physical atom is thus
represented by $P$ replicas, or beads, joined by harmonic springs. The
physical potential is evaluated independently on every bead, while the
springs encode the quantum delocalization of the nuclei in imaginary time.
When $P=1$, Eq.~\eqref{eq:ring-polymer-hamiltonian} reduces to the classical
configurational distribution. Larger $P$ is required as temperature falls or
as the maximum physical vibrational frequency increases; convergence must be
checked for the chosen isotope masses and temperature range.

Two useful geometrical quantities are the bead centroid and radius of
gyration,
\begin{equation}
  \bar{\bm R}_i=\frac{1}{P}\sum_{s=1}^{P}\bm R_i^{(s)},
  \qquad
  r_{g,i}^2=\frac{1}{P}\sum_{s=1}^{P}
  |\bm R_i^{(s)}-\bar{\bm R}_i|^2.
  \label{eq:ring-polymer-centroid}
\end{equation}
The centroid describes the mean imaginary-time position, while the spread of
the beads measures nuclear quantum delocalization. For a position-dependent
operator $A(\bm R)$, the primitive estimator has the form
\begin{equation}
  \langle A\rangle
  =\lim_{P\to\infty}
   \left\langle\frac{1}{P}\sum_{s=1}^{P}
   A(\bm R^{(s)})\right\rangle_{H_P}.
  \label{eq:path-integral-estimator}
\end{equation}
Kinetic energy, pressure and heat capacity require specialized estimators or
thermodynamic derivatives; they are not obtained by treating the fictitious
ring-polymer kinetic energy as the physical nuclear kinetic energy.

\subsubsection{Sampling the ring polymer}

PIMD introduces fictitious momenta in order to sample the ring-polymer
distribution using MD algorithms. A discrete Fourier or normal-mode
transformation diagonalizes the free ring-polymer spring term. Its mode
frequencies are
\begin{equation}
  \omega_k=2\omega_P\sin\left(\frac{k\pi}{P}\right),
  \qquad k=0,\ldots,P-1,
  \label{eq:ring-polymer-modes}
\end{equation}
where $k=0$ is the centroid mode. The large separation between the slow
physical motion and the stiff internal modes makes naive sampling inefficient.
Normal-mode or staging coordinates, mass preconditioning, multiple time
steps, and mode-specific thermostats address this stiffness. The PILE
thermostat applies analytically tuned Langevin friction to the known free
ring-polymer modes, while the PILE-L variant also thermostats the centroid
\cite{pile}.

The standard primitive factorization has an asymptotic discretization error
of order $\mathcal O(P^{-2})$. The fourth-order Suzuki--Chin factorization
used in the reference work replaces the bead potential by a bead-dependent
effective potential of the general form
\begin{equation}
  \tilde U_s(\bm R;\beta)
  =w_s U(\bm R)
   +w_s d_s\left(\frac{\beta\hbar}{P}\right)^2
    \sum_i\frac{|\bm F_i(\bm R)|^2}{m_i},
  \label{eq:suzuki-chin}
\end{equation}
where $w_s$ and $d_s$ are factorization coefficients that alternate between
beads. The force-squared correction accounts for the leading non-commutation
of kinetic and potential operators and gives fourth-order convergence for
appropriate thermodynamic observables. Direct differentiation would involve
Hessians; finite-difference implementations make the scheme practical with
general force providers \cite{highorder}. Ring-polymer contraction can further
reduce cost by evaluating slowly varying, long-range components of the
potential on fewer contracted beads, while rapidly varying short-range terms
remain evaluated on the full polymer.

\subsubsection{What PIMD does and does not predict}

The trajectories generated by ordinary PIMD are a sampling device for static
quantum equilibrium averages. The bead index represents imaginary time, and
the artificial bead masses, thermostats and ring-polymer vibrations have no
general identification with real nuclear motion. Consequently, ordinary PIMD
does not directly yield quantum time-correlation functions, diffusion
coefficients or vibrational spectra. Ring-polymer molecular dynamics (RPMD)
and centroid molecular dynamics are related but additional approximations for
real-time correlation functions; they should not be conflated with PIMD.

PIMD captures nuclear zero-point fluctuations, tunnelling contributions to
equilibrium statistics, and anharmonic sampling on the supplied PES, subject
to convergence in bead count, timestep, sampling length and system size. It
does not correct errors in that PES, include electronic excited states, or
automatically change an NVT calculation into a constant-pressure one. For the
heat capacity in the reference study, the ring polymer was sampled in a
constant-pressure ensemble and \(\cpheat\) was obtained from converged
enthalpy derivatives at neighboring temperatures, as discussed below.

\subsection{Thermodynamic definitions}

\subsubsection{Constant-volume and constant-pressure heat capacities}

For a system at equilibrium, the constant-volume and constant-pressure heat
capacities are
\begin{align}
  C_V &= \left(\frac{\partial \langle E\rangle}{\partial T}\right)_{V,N},
  \\
  C_P &= \left(\frac{\partial \langle H\rangle}{\partial T}\right)_{P,N},
  \qquad H=E+PV,
\end{align}
where $E$ is the total internal energy and $H$ is the enthalpy. For a stable
single phase,
\begin{equation}
  C_P-C_V = \frac{TV\alpha_V^2}{\kappa_T},
  \label{eq:cp-cv}
\end{equation}
where $\alpha_V=V^{-1}(\partial V/\partial T)_P$ is the volumetric thermal
expansion coefficient and
$\kappa_T=-V^{-1}(\partial V/\partial P)_T$ is the isothermal compressibility.
Equation~\eqref{eq:cp-cv} explains why \cv{} and \cpheat{} need not be
interchangeable even when their numerical difference is small for a
particular solid.

In an equilibrated classical canonical ensemble, \cv{} can also be obtained
from total-energy fluctuations,
\begin{equation}
  C_V = \frac{\langle E^2\rangle-\langle E\rangle^2}
  {k_{\mathrm B}T^2}.
  \label{eq:fluctuation-cv}
\end{equation}
The analogous isothermal--isobaric expression uses enthalpy fluctuations.
These relations require a stationary ensemble, long sampling, and uncertainty
estimation that accounts for time correlation. Moreover, classical nuclear
dynamics gives equipartition rather than the quantum freezing of
high-frequency modes. The present workflow obtains its classical
constant-pressure term by differentiating replica-averaged NPT enthalpies on a
temperature grid. It obtains the quantum correction separately from harmonic
normal-mode frequencies.

\subsubsection{Why classical and quantum predictions differ}

In the classical harmonic limit, every quadratic vibrational degree of
freedom contributes $k_{\mathrm B}$ to \cv, producing the Dulong--Petit limit
at all temperatures. Quantum mechanically, a mode of angular frequency
$\omega$ becomes thermally populated only when $\hbar\omega$ is comparable to
or smaller than $k_{\mathrm B}T$. High-frequency framework and molecular modes
are therefore partially frozen near room temperature. Classical MD can still
describe anharmonic structural motion, but it cannot generate the correct
quantum populations simply by extending the trajectory.

\section{Reference study: Kapil \emph{et al.}}

Kapil \emph{et al.}\ studied empty MOF-5 and cells containing 50, 100, and 150
methane molecules using classical MD, Suzuki--Chin path-integral molecular
dynamics (PIMD), and quantum harmonic normal-mode analysis
\cite{kapil2019}. Their work provides both the physical motivation and an
important limit on what the present calculation can claim.

\subsection{Force field and starting configurations}

The reference work used analytical QuickFF force fields derived from
electronic-structure calculations, not PET-MAD \cite{quickff}. Framework
covalent terms were
fitted to B3LYP cluster calculations, with MBIS charges obtained from a
periodic PBE density. Nonbonded interactions combined Gaussian-distributed
electrostatics with an MM3 Buckingham model and a \SI{15}{\angstrom} cutoff
with tail corrections. Methane was inserted with RASPA, configurations with
unrealistic contacts were rejected, and canonical Monte Carlo was used before
dynamics. Consequently, agreement of timestep and temperature alone would not
make the present PET-MAD simulations a numerical reproduction of the paper.

\subsection{Published simulation protocols}

\begin{table}[htbp]
\centering
\small
\begin{tabular}{@{}p{0.25\linewidth}p{0.34\linewidth}p{0.34\linewidth}@{}}
\toprule
Parameter & Classical MD & Suzuki--Chin PIMD \\
\midrule
State points & 100--\SI{500}{K}, \SI{1}{bar} & 100--\SI{500}{K}, \SI{1}{bar} \\
Cell dynamics & Fully flexible cell, zero deviatoric stress & Fluctuating volume with cubic cell shape retained \\
Integrator & Verlet, \SI{0.5}{fs} & BAOAB-type multiple-time-step scheme \\
Thermostat & Three-element Nos\'e--Hoover chain, \SI{100}{fs} relaxation & PILE-L ring-polymer thermostat; \SI{100}{fs} relaxation \\
Barostat & Martyna--Tobias--Klein, \SI{1000}{fs} relaxation & Bussi--Zykova--Parrinello form adapted to Suzuki--Chin; \SI{1000}{fs} relaxation \\
Quantum discretization & Not applicable & 64 beads; long-range forces contracted to 8 beads \\
Force timesteps & \SI{0.5}{fs} & \SI{0.25}{fs} short range; \SI{1}{fs} long range \\
Empty MOF-5 sampling & One \SI{500}{ps} trajectory & Five independent \SI{125}{ps} trajectories \\
Loaded sampling & Five independent \SI{500}{ps} trajectories & Thirty independent \SI{50}{ps} trajectories \\
Discarded equilibration & First \SI{100}{ps} & First \SI{25}{ps} \\
Heat-capacity estimator & Classical thermodynamic sampling & Centered enthalpy derivative with $\Delta T=\SI{25}{K}$ \\
\bottomrule
\end{tabular}
\caption{Principal simulation parameters reported by Kapil \emph{et al.}\ Numerical random seeds and output strides were not reported.}
\label{tab:kapil-methods}
\end{table}

The associated Nos\'e--Hoover-chain, constant-pressure, path-integral, and
PILE-L methods are described in Refs.~\cite{nhc,mtk,ipi,pile}. Their production
\cpheat{} curves were evaluated from neighboring enthalpies,
\begin{equation}
  C_P(T) \simeq
  \frac{\mathcal H(T+\Delta T)-\mathcal H(T-\Delta T)}{2\Delta T},
  \qquad \Delta T=\SI{25}{K},
  \label{eq:finite-difference-cp}
\end{equation}
rather than from a single short trajectory. Every temperature point therefore
depends on converged, statistically independent simulations at adjacent
temperatures.

For the harmonic comparison, optimized structures were used to construct
Cartesian Hessians with Yaff and normal modes with TAMkin. Empty MOF-5 was
tested with both QuickFF and UFF4MOF. The loaded 100-methane calculation used
five different Monte Carlo configurations, each optimized into a local
minimum before its harmonic spectrum was evaluated.

\subsection{Physical conclusions relevant to this project}

The main findings of the reference study are:

\begin{itemize}
  \item Empty MOF-5 is quantum mechanical but nearly harmonic over the studied
        range. Its quantum harmonic \cv{} is close to the PIMD \cpheat{}, with
        a room-temperature gravimetric value of approximately
        \SI{0.7}{J\per g\per K}.
  \item Classical MD overestimates the empty-framework heat capacity because
        it populates all vibrational modes according to equipartition.
  \item Methane loading introduces strongly anharmonic host--guest motion. The
        PIMD curves develop a loading-dependent minimum around \SI{200}{K},
        caused mainly by changes in attractive host--guest interactions as
        confined methane becomes more mobile.
  \item A harmonic calculation cannot represent diffusion, cage-to-cage
        motion, or the temperature-dependent localization of methane. Nuclear
        quantum effects and anharmonic sampling are both needed for a
        quantitative loaded-system \cpheat{}.
  \item The approximate combination
        \begin{equation}
          C \approx C_{\mathrm{qn}}^{\mathrm{har}}
          - C_{\mathrm{cl}}^{\mathrm{har}}
          + C_{\mathrm{cl}}^{\mathrm{anh}}
          \label{eq:quantum-correction}
        \end{equation}
        qualitatively recovers the loaded minimum and becomes useful above
        roughly \SI{200}{K} in their 100-methane test. This is the approximation
        implemented in the current repository.
\end{itemize}

\section{Present computational methodology}

\subsection{Structure preparation and methane insertion}

The repository host structure is the conventional cubic MOF-5 cell with 424
atoms and a cell parameter of approximately \SI{25.866}{\angstrom}. The loaded
campaign is generated by
\texttt{python -m mof\_heat\_capacity.protocols.loaded}. For every temperature and
replica, methane molecules are centered, uniformly randomly rotated using a
quaternion, and placed at uniformly sampled fractional coordinates. Periodic
atom--atom distances are checked against all previously placed atoms. A
100-methane structure contains 924 atoms. Independent insertion and velocity
seeds prevent the nominal replicas from sharing an identical initial state.

This procedure prevents severe initial overlaps, but it is not an adsorption
equilibrium calculation. In particular, it does not reproduce the RASPA
insertion and Monte Carlo equilibration used by Kapil \emph{et al.} Several
independent insertion seeds and adequate subsequent equilibration are required
before drawing conclusions about loading-dependent thermodynamics.

\subsection{PET model and implementation}
\label{sec:pet-implementation}

Section~\ref{sec:pet-architecture} describes the original PET architecture.
The local PET-MAD and PET-SOL conversion metadata both specify the settings
in Table~\ref{tab:pet-settings}. These describe the supplied checkpoints,
not every PET model; identical architecture settings do not imply identical
weights or training data. In particular, the original paper's default of
three attention layers per block differs from the one layer used here.

\begin{table}[htbp]
\centering
\caption{Architecture settings recorded in the local PET-MAD and PET-SOL
PET-JAX metadata. The species count is embedding capacity, not a measure of
training coverage.}
\label{tab:pet-settings}
\begin{tabular}{lr}
\toprule
Setting & Value \\
\midrule
Outer cutoff & \SI{8.0}{\angstrom} \\
Cutoff switching width & \SI{0.5}{\angstrom} \\
Adaptive neighbor-selection setting & 40 \\
Message-passing blocks & 3 \\
Attention layers per block & 1 \\
Token width, $d_{\mathrm{PET}}$ & 256 \\
Attention heads & 8 \\
Transformer feed-forward width & 512 \\
Node feature width, \texttt{d\_node} & 1024 \\
Readout feature width, \texttt{d\_head} & 256 \\
Species-embedding capacity & 102 \\
\bottomrule
\end{tabular}
\end{table}

The locally inspected PyTorch implementation (\texttt{metatrain} 2026.3.1,
\texttt{pet/modules/transformer.py}) embeds four geometric inputs,
$(x_{ij},y_{ij},z_{ij},r_{ij})$, rather than only the three Cartesian
components in the original Eq.~\eqref{eq:pet-token}. Initial edge messages
are species embeddings. Wider central features are contracted to the token
width for attention and expanded afterward. With the dimensions in
Table~\ref{tab:pet-settings}, each attention head has width $256/8=32$;
\texttt{d\_head=256} denotes a readout width. Padded neighbors are masked,
and logarithms of cutoff factors are added to attention logits, implementing
the weighting in Eq.~\eqref{eq:self-attention} up to numerical regularization.
In the residual feature path of \texttt{pet/model.py}, a reverse-neighbor
index routes outgoing tokens to receiving environments; averaging the old
message with the reversed new output updates the edge state. These are
implementation details of the installed PyTorch package; the checkpoint
version and PET-JAX conversion must remain compatible with them.

The model weights remain fixed at inference. The repository delegates the
network implementation to the PET software stack: its export helper calls
\texttt{upet.save\_upet} to turn the selected checkpoint into a portable
metatomic model, which LAMMPS evaluates during production MD. Harmonic
analysis converts the matching checkpoint into PET-JAX parameter and metadata
files, then uses SADMOF to differentiate its scalar energy. The setup script
pins PET-JAX to commit \texttt{ed0add4}; conversion transfers parameters rather
than retraining a potential. Species baselines have zero Cartesian derivatives
at fixed composition, whereas the learned energy scale must be retained in
forces and Hessians.

Using matching checkpoints, exports, and conversions is essential to keep
MD and harmonic analysis on a consistent PES. This must be checked by
comparing energies and forces, including adaptive-cutoff derivative
conventions, rather than inferred from matching architecture settings alone.
The sparse Hessian path stops adaptive-cutoff derivatives, as described below.
The selected metatomic export must expose stress, and campaign configurations
record stress validation explicitly rather than inferring capability from a
model filename. PET is the architecture; PET-MAD is a general-purpose fitted
potential, whose energy, force, stress, and curvature accuracy for MOF-5 and
its methane guests must still be validated \cite{petmad}.

\subsection{Classical NPT molecular dynamics}

The statistical and dynamical foundations of classical MD are summarized in
Section~\ref{sec:classical-md-theory}. Only methane-loaded MOF-5 is propagated.
LAMMPS evaluates a metatomic PET model in a flexible triclinic cell at an
isotropic target pressure of \SI{1}{bar}. Temperature and pressure are
controlled by a Nos\'e--Hoover-chain/MTTK NPT fix, with chain length 3,
thermostat damping time \SI{100}{fs}, and barostat damping time
\SI{1000}{fs}. The production template specifies $10^6$ steps of
\SI{0.5}{fs}, hence \SI{500}{ps}, saves data every 100 steps
(\SI{0.05}{ps}), and discards the first \SI{100}{ps}. The initial campaign
uses one replica at 100, 200, 300, 400, and \SI{500}{K}. This coarse
\SI{100}{K} spacing reduces cost but limits resolution of curvature in the
enthalpy curve and of the heat-capacity minimum; denser temperatures and
independent replicas remain convergence extensions.

The exported model must expose stress and must be validated for the intended
MOF/loading states before NPT use. Mean loaded-system enthalpies are computed
from the stationary portions of the trajectories,
\begin{equation}
  \langle H\rangle_T
  =\left\langle E+P_{\mathrm{ext}}V\right\rangle_T,
\end{equation}
then averaged across replicas and differentiated with respect to temperature.
This classical contribution retains guest diffusion, localization changes,
and other anharmonic effects sampled by the MLIP.

\subsection{Harmonic approximation and Cartesian Hessian}

Let $\bm R_0$ be a stationary reference structure and
$\bm u=\bm R-\bm R_0$. Expanding the potential energy gives
\begin{equation}
  U(\bm R_0+\bm u)=U(\bm R_0)
  +\underbrace{\nabla U(\bm R_0)\cdot\bm u}_{=0\ \mathrm{at\ a\ minimum}}
  +\frac{1}{2}\bm u^{\mathrm T}\bm H\bm u+\mathcal O(u^3),
  \label{eq:taylor}
\end{equation}
with Cartesian Hessian
\begin{equation}
  H_{i\alpha,j\beta}
  =\frac{\partial^2 U}{\partial R_{i\alpha}\partial R_{j\beta}}
  =-\frac{\partial F_{i\alpha}}{\partial R_{j\beta}}.
  \label{eq:hessian}
\end{equation}
For $N$ atoms, $\bm H$ has dimensions $3N\times3N$; for the 100-methane cell
this is $2772\times2772$.

PET-JAX differentiates the model computation graph rather than using displaced
finite differences. The implementation uses forward-over-reverse
Hessian--vector products,
\begin{equation}
  \bm H\bm v=
  \left.\frac{\mathrm d}{\mathrm d\epsilon}
  \nabla U(\bm R+\epsilon\bm v)\right|_{\epsilon=0}.
\end{equation}
Because a local message-passing model couples only graph-connected atoms, the
Hessian has a block-sparse structure. SADMOF \cite{sadMof} constructs a
graph-derived sparsity pattern and ASDEX colors its columns so that several independent
columns can be reconstructed from one differentiated evaluation. The current
settings are single precision, graph depth \texttt{hops=3}, HVP chunk size 1,
and rematerialization enabled. The implementation calls
\texttt{get\_energy\_fn(..., no\_shadow=True)}, which stops derivatives
through the adaptive cutoff and confines the sparse treatment to the selected
graph. Thus this is a specified derivative convention; unrestricted AD
through adaptive selection includes additional couplings. Its consistency
with the MD force convention requires explicit comparison.

These settings are not established convergence values. In particular,
\texttt{hops} is the depth of the Hessian sparsity pattern, not the number of
PET message-passing blocks. The local SADMOF implementation identifies seven
hops as the structurally exact pattern for its three-layer PET-MAD-S graph
construction \cite{sadMof}. The project's three-hop setting is therefore a
truncation: it must be enlarged until frequencies and \cv{} stop changing,
and checked against a dense or less aggressively truncated reference where
feasible. Precision and the adaptive-cutoff derivative convention also need
validation, especially for the low-frequency modes.

\subsection{Normal modes and quantum harmonic heat capacity}

The Hessian is symmetrized and mass weighted to obtain
\begin{equation}
  D_{i\alpha,j\beta}
  =\frac{H_{i\alpha,j\beta}}{\sqrt{m_i m_j}},
  \qquad
  \bm D\bm e_k=\omega_k^2\bm e_k.
\end{equation}
The code enforces the force-constant acoustic sum rule before diagonalization,
\begin{equation}
  \sum_j H_{i\alpha,j\beta}=0,
\end{equation}
so that uniform translations have zero restoring force. The upstream SADMOF
observable maps negative eigenvalues to zero frequency for compatibility with
its reference heat-capacity convention. The present workflow instead
diagonalizes the mass-weighted Hessian locally and retains unstable eigenvalues
as signed negative frequencies. Hybrid assembly rejects unexpected imaginary
modes and permits at most the three translational near-zero modes, so a
thermally displaced frame or saddle point cannot silently enter the correction.

For each retained mode, define
\begin{equation}
  x_k=\frac{\hbar\omega_k}{k_{\mathrm B}T}
      =\frac{hc\tilde\nu_k}{k_{\mathrm B}T},
\end{equation}
where $\tilde\nu_k$ is the wavenumber. The quantum harmonic contribution is
\begin{equation}
  C_{V,k}=k_{\mathrm B}
  \left[\frac{x_k/2}{\sinh(x_k/2)}\right]^2,
\end{equation}
and the gravimetric result is
\begin{equation}
  c_V(T)=\frac{N_{\mathrm A}}{M_{\mathrm{cell}}}
  \sum_{k\in\mathcal V} C_{V,k},
  \label{eq:harmonic-cv}
\end{equation}
in \si{J\per g\per K}. The mass $M_{\mathrm{cell}}$ includes methane for a
loaded system, so both the spectrum and the normalization change with loading.

One Hessian produces the entire configured 100--\SI{500}{K} curve in
\SI{10}{K} increments. The temperature appearing in
Eq.~\eqref{eq:harmonic-cv} controls oscillator populations; it does not require
a separate MD trajectory at every analysis temperature when the force
constants are assumed temperature independent.

\subsection{Optimized minima and the hybrid correction}

The empty-framework calculation starts from the equilibrated MOF-5 structure
already available to the project. Its coordinates are relaxed at fixed cell
with the selected MLIP and one AD Hessian is evaluated. No empty-MOF MD is
needed merely to obtain a Hessian reference.

For the loaded system, final structures from several independent classical MD
replicas are quenched at fixed cell with the same MLIP. A Hessian is accepted
only after the relaxation reaches the chosen force tolerance and its signed
spectrum is checked for unexpected imaginary modes. Raw finite-temperature
frames are not treated as normal-mode reference structures. If $N_m$
independent optimized minima are retained, the mean harmonic quantum
correction is
\begin{equation}
  \overline{\Delta c^{\mathrm{har}}}(T)
  =\frac{1}{N_m}\sum_{s=1}^{N_m}
  \left[c_{\mathrm{qn},s}^{\mathrm{har}}(T)
       -c_{\mathrm{cl},s}^{\mathrm{har}}\right],
  \label{eq:minimum-average}
\end{equation}
with sample variance
\begin{equation}
  s_{\Delta c}^2(T)=\frac{1}{N_m-1}
  \sum_{s=1}^{N_m}\left[
    \Delta c_s^{\mathrm{har}}(T)
    -\overline{\Delta c^{\mathrm{har}}}(T)\right]^2.
\end{equation}

For each minimum, the quantum and classical harmonic terms use exactly the
same retained-mode mask. The classical term is $k_{\mathrm B}$ per retained
mode. The complete loaded-system spectrum is used because framework and
methane motion can mix; the empty Hessian alone omits methane quantization and
cannot supply this correction.

The final estimator is
\begin{equation}
  c_P^{\mathrm{approx}}(T)
  =c_{P,\mathrm{cl}}^{\mathrm{anh}}(T)
  +\overline{\Delta c^{\mathrm{har}}}(T),
  \qquad
  c_{P,\mathrm{cl}}^{\mathrm{anh}}(T)
  =\frac{\mathrm d\langle H\rangle_{\mathrm{cl}}}{\mathrm dT}.
  \label{eq:hybrid-cp}
\end{equation}
Extensive terms are combined before division by the total mass of the loaded
cell. Replica-aware resampling propagates uncertainty in the enthalpy curve;
the spread across optimized minima measures sensitivity of the harmonic
correction. The Hessian contribution is formally a \cv{} correction added to a
classical \cpheat{} term, following the small solid-state $C_P-C_V$ assumption
used in the reference approximation. This mixing and the approximation's
limited validation below roughly \SI{200}{K} must be stated with the result.

\section{Results and discussion}
\label{sec:results}

The earlier trajectory and raw-frame Hessian calculations were exploratory and
are not results of the focused workflow defined here. Their numerical tables
and generated overview figure have therefore been removed from this report.
Production results should be added only after the new loaded replicas,
optimized-minimum Hessians, and hybrid analysis have completed and passed the
checks below.

The primary reported result will be a temperature-resolved table or curve
containing the replica-averaged classical enthalpy derivative, the mean loaded
harmonic quantum correction, their separate uncertainties, and the combined
$c_P^{\mathrm{approx}}$. The empty quantum-harmonic curve will be shown as a
separately normalized reference. Density, cell response, framework stability,
methane mobility, imaginary-mode counts, and relaxation residual forces will
support interpretation rather than replace the heat-capacity result.

\section{Recommended production and analysis protocol}

\subsection{Empty MOF-5 reference}

The equilibrated empty-framework structure is used directly; an empty-MOF MD
campaign would add cost without supplying the anharmonic loaded-system term
needed by Eq.~\eqref{eq:hybrid-cp}.

\begin{enumerate}
  \item Relax the atomic coordinates with the same PET-MAD checkpoint used for
        the Hessian. Report maximum residual force and do not relax the cell
        unless a stress-capable, validated model is adopted.
  \item Check stability using signed eigenvalues and inspect low-frequency
        modes. A large number of unstable modes invalidates the harmonic
        interpretation.
  \item Converge float32 against float64, graph hops, HVP chunking, acoustic
        sum-rule treatment, and frequency threshold.
  \item Converge the conventional-cell result with respect to supercell size or
        phonon wave-vector sampling. A single periodic cell Hessian samples
        only the corresponding finite set of modes.
  \item Retain the resulting curve as a separately normalized reference and
        compare its room-temperature value with the approximately
        \SI{0.7}{J\per g\per K} reference scale, while stating that PET-MAD
        replaces QuickFF and that model error remains to be quantified.
\end{enumerate}

\subsection{Loaded MOF-5 production and Hessians}

For the 100-methane system, the production sequence is:

\begin{enumerate}
  \item Prepare independent methane placements and velocity seeds for each
        replica on the enthalpy temperature grid.
  \item Run flexible-cell classical NPT only for the loaded system. Determine
        equilibration from enthalpy, cell, framework, and methane-mobility
        diagnostics rather than assuming the template duration is sufficient.
  \item Average enthalpy over the stationary part of each replica, combine
        replica uncertainties, and differentiate the temperature dependence.
  \item Select representative independent loaded configurations, quench each
        one to a fixed-cell minimum with the same MLIP, and reject unconverged
        or unstable minima.
  \item Compute one AD Hessian per accepted minimum. Converge precision, graph
        hops, HVP chunking, acoustic treatment, frequency threshold, and cell
        size on the quantum-minus-classical correction.
  \item Average the loaded corrections, combine them with the classical
        enthalpy derivative using Eq.~\eqref{eq:hybrid-cp}, and report MD and
        minimum-to-minimum uncertainties separately as well as jointly.
\end{enumerate}

This calculation is intentionally not an exact reproduction of the reference
study. A small number of path-integral state points may be considered later as
validation if quantitative accuracy below approximately \SI{200}{K} is
required, but they are not part of the production workflow.

\subsection{Equilibration and convergence criteria}

Production decisions should be based on observables, not a fixed waiting time
alone. At minimum, inspect:

\begin{itemize}
  \item running and block averages of temperature, potential energy, and total
        energy;
  \item maximum and distribution of force magnitudes;
  \item framework bond integrity, pore geometry, and cell consistency;
  \item methane center-of-mass distributions, host--guest distances, and
        mean-squared displacements;
  \item sensitivity to timestep and thermostat friction;
  \item agreement among independent seeds;
  \item convergence of the enthalpy derivative with replica count, sampling
        length, and temperature spacing; and
  \item convergence of the harmonic correction with independent minima,
        precision, sparsity depth, cell size, and mode treatment.
\end{itemize}

The timestep should first be validated in short NVE tests by checking energy
drift. Thermostat convergence should then be tested separately, because a
thermostat can conceal an integration error while maintaining the desired
kinetic temperature.

\section{Implementation, outputs, and reproducibility}

The production campaign is prepared by
\texttt{python -m mof\_heat\_capacity.protocols.loaded} and submitted by
\texttt{simulation/submit\_loaded\_md.sh}. These commands create independent
loaded structures and LAMMPS/metatomic NPT runs while deliberately rejecting a
zero-methane MD campaign. Trajectory diagnostics are submitted with
\texttt{properties/submit\_analysis.sh}.

The command \texttt{properties/submit\_heat\_capacity.sh} takes final structures
from selected loaded replicas, relaxes them with
\texttt{mof\_heat\_capacity.structures.relax}, and passes the optimized
structures to \texttt{mof\_heat\_capacity.analysis.harmonic}. The same command
relaxes and evaluates the supplied empty reference directly. Harmonic archives
store the run and model identity, optimized structure, temperature grid,
frequency array, gravimetric \cv{} curve, precision, graph hops, HVP chunk
size, and rematerialization choice. Finally,
\texttt{properties/submit\_hybrid\_analysis.sh} differentiates the loaded NPT
enthalpies and writes the classical term, harmonic correction, combined curve,
uncertainties, and provenance as NPZ, CSV, and JSON files.

Cluster calculations use one node, one task, and one GPU because neither the MD
nor Hessian code is MPI or multi-GPU parallel. Independent temperature,
replica, or optimized-minimum jobs are the available level of parallelism.
Slurm job identifiers and resource use should be retained; elapsed time and peak
memory can be recovered with, for example,
\begin{verbatim}
sacct -j JOB_ID --format=JobID,State,Elapsed,MaxRSS,AllocTRES
\end{verbatim}
Wall time and memory requests should be based on representative jobs.
Separate output directories and
prefixes are required for every loading, temperature, seed, and replica.

A reproducible result should record:

\begin{itemize}
  \item source structure, cell, atom count, chemical composition, loading, and
        methane-insertion seed and rejection criteria;
  \item PET-MAD checkpoint hash, metatomic export, PET-JAX conversion, software
        environment, device, and numerical precision;
  \item ensemble, target temperature, timestep, thermostat and friction,
        number of steps, output stride, velocity seed, equilibration interval,
        production interval, and restart history;
  \item Hessian source replica, relaxation state, residual forces, graph hops,
        HVP chunk size, rematerialization, acoustic sum rule, frequency
        threshold, and temperature grid; and
  \item independent-replica count, autocorrelation or block length,
        uncertainty estimator, and all convergence tests.
\end{itemize}

\section{Limitations and interpretation}

The principal limitations of the planned result are:

\begin{itemize}
  \item PET-MAD and QuickFF describe different potential-energy surfaces;
        agreement with the paper cannot be assumed.
  \item Flexible-cell NPT requires a stress-capable MLIP and a MOF-5-specific
        stress benchmark; an output capability flag alone is not validation.
  \item Classical nuclei are used in MD; nuclear quantum effects enter only
        through the harmonic oscillator formula.
  \item The random insertion is not an adsorption or Monte Carlo equilibrium
        distribution.
  \item A finite number of quenched structures may not represent every
        thermodynamically important loaded basin. Their spread measures
        sensitivity but is not automatically a free-energy-weighted average.
  \item The harmonic correction cannot itself represent methane diffusion or
        barrier crossing; those effects enter only through classical MD.
  \item Adding a harmonic \cv{} correction to a classical \cpheat{} term assumes
        that the corresponding pressure--volume quantum correction is small.
  \item The approximation was quantitatively successful in the reference
        100-methane system mainly above about \SI{200}{K}; transfer to this
        MLIP, other loadings, and lower temperatures requires validation.
  \item Finite-size, sparse-Hessian depth, precision, temperature-derivative,
        model-validity, and independent-replica errors must all be quantified.
\end{itemize}

These qualifications define the boundary between an efficient, physically
motivated approximation and an exact quantum-anharmonic thermodynamic result.
They should be carried into every plot, table, and comparison with experiment.

\section{Conclusions}

The focused workflow avoids both unnecessary empty-framework MD and the cost
of a full production path-integral campaign. Classical NPT trajectories are
reserved for loaded MOF-5, where they are needed to describe anharmonic guest
and host--guest motion. Sparse AD Hessians are evaluated only at converged
fixed-cell minima and provide the loaded-system quantum-minus-classical
harmonic correction. The equilibrated empty structure contributes one direct
reference Hessian.

The final quantity is the classical loaded-system enthalpy derivative plus the
mean loaded harmonic correction. Its credibility depends on stress validation,
stationary and replicated MD, stable optimized minima, numerical Hessian
convergence, consistent mode and mass bookkeeping, and explicit uncertainty
propagation. This strategy implements the cost-saving approximation motivated
by Kapil \emph{et al.} without attempting to reproduce every calculation in
their study.

\begin{thebibliography}{11}

\bibitem{kapil2019}
V. Kapil \emph{et al.},
\emph{Modeling the Structural and Thermal Properties of Loaded
Metal--Organic Frameworks: An Interplay of Quantum and Anharmonic
Fluctuations}, Journal of Chemical Theory and Computation \textbf{15},
3237--3249 (2019),
\href{https://doi.org/10.1021/acs.jctc.8b01297}{doi:10.1021/acs.jctc.8b01297}.

\bibitem{ase}
A. H. Larsen \emph{et al.},
\emph{The Atomic Simulation Environment---a Python Library for Working with
Atoms}, Journal of Physics: Condensed Matter \textbf{29}, 273002 (2017),
\href{https://doi.org/10.1088/1361-648X/aa680e}{doi:10.1088/1361-648X/aa680e}.

\bibitem{petmad}
A. Mazitov \emph{et al.},
\emph{PET-MAD as a Lightweight Universal Interatomic Potential for Advanced
Materials Modeling}, Nature Communications \textbf{16}, 10653 (2025),
\href{https://doi.org/10.1038/s41467-025-65662-7}{doi:10.1038/s41467-025-65662-7}.

\bibitem{pet}
S. N. Pozdnyakov and M. Ceriotti,
\emph{Smooth, Exact Rotational Symmetrization for Deep Learning on Point
Clouds}, Advances in Neural Information Processing Systems \textbf{36}
(2023),
\href{https://proceedings.neurips.cc/paper_files/paper/2023/hash/fb4a7e3522363907b26a86cc5be627ac-Abstract-Conference.html}{conference paper};
the supplied article is
\href{https://arxiv.org/abs/2305.19302v3}{arXiv:2305.19302v3}
(6 February 2024), including Appendices A--H.

\bibitem{quickff}
L. Vanduyfhuys \emph{et al.},
\emph{QuickFF: A Program for a Quick and Easy Derivation of Force Fields for
Metal--Organic Frameworks from Ab Initio Input}, Journal of Computational
Chemistry \textbf{36}, 1015--1027 (2015).

\bibitem{ipi}
M. Ceriotti, J. More, and D. E. Manolopoulos,
\emph{i-PI: A Python Interface for Ab Initio Path Integral Molecular Dynamics
Simulations}, Computer Physics Communications \textbf{185}, 1019--1026 (2014).

\bibitem{mtk}
G. J. Martyna, D. J. Tobias, and M. L. Klein,
\emph{Constant Pressure Molecular Dynamics Algorithms}, Journal of Chemical
Physics \textbf{101}, 4177--4189 (1994).

\bibitem{nhc}
G. J. Martyna, M. L. Klein, and M. Tuckerman,
\emph{Nos\'e--Hoover Chains: The Canonical Ensemble via Continuous Dynamics},
Journal of Chemical Physics \textbf{97}, 2635--2643 (1992).

\bibitem{pile}
M. Ceriotti, M. Parrinello, T. E. Markland, and D. E. Manolopoulos,
\emph{Efficient Stochastic Thermostatting of Path Integral Molecular
Dynamics}, Journal of Chemical Physics \textbf{133}, 124104 (2010).

\bibitem{highorder}
V. Kapil, J. Behler, and M. Ceriotti,
\emph{High Order Path Integrals Made Easy}, Journal of Chemical Physics
\textbf{145}, 234103 (2016),
\href{https://doi.org/10.1063/1.4971438}{doi:10.1063/1.4971438}.

\bibitem{sadMof}
M. F. Langer, M. Ceriotti, and collaborators,
\emph{SADMOF: Sparse Automatic Differentiation for Heat Capacity of
Metal--Organic Frameworks}, project software and methodology under
development (version used by this repository, 2026).

\end{thebibliography}

\end{document}
